%% toc: P3 36 36.6
\section{Proofs for partially ordered quantifier prefixes}
\label{sec:annex-p03-w03-p03-044-po-prefixes}

Throughout, \(s[a/z]\) is the assignment that agrees with \(s\) except that
it sends \(z\) to \(a\).  Structures have nonempty universes.

\begin{annexitem}{annex:p03-w03-p03-044-linear-skolem}{thm:p03-w03-p03-044-linear-skolem}{Linear prefixes and witness functions}
\begin{proof}
For a linear prefix \(\pi=\mathsf Q_1z_1\cdots\mathsf Q_nz_n\), the
existential variable \(z_j\) has as argument list the universal variables
\(z_i\) with \(i<j\), in increasing order.  Call a family of functions
\((f_{z_j})_{z_j\in E_\pi}\) of these arities a \emph{witness family for
\(\pi\varphi\) at \(s\)} when, for every assignment of values to the
universal variables of \(\pi\), the formula \(\varphi\) holds at the
assignment that extends \(s\) by those values and gives each existential
\(z_j\) the value of \(f_{z_j}\) at the values of its argument list.  By
\cref{def:p03-w03-p03-044-henkin-prefix}, \(M\entails Q_\pi\varphi\,[s]\)
holds iff a witness family exists.  Both parts are proved together by
induction on \(n\), for all formulas \(\varphi\) and all assignments \(s\)
simultaneously; part~(2) uses choice only in the universal case.

\emph{Base \(n=0\).}  The triple is
\((\varnothing,\varnothing,\varnothing)\).  The only witness family is
empty, and the clause requires exactly \(M\entails\varphi\,[s]\).  This is
also the ordinary meaning of the empty prefix.

\emph{Step.}  Let \(\pi=\mathsf Q_1z_1\pi'\) with
\(\pi'=\mathsf Q_2z_2\cdots\mathsf Q_nz_n\), and assume both parts for
\(\pi'\), every formula and every assignment.

\emph{Case \(\mathsf Q_1=\forall\).}  The variable \(z_1\) precedes every
existential variable, so \(D_\pi=D_{\pi'}\cup\{(z_1,z_j):z_j\in E_\pi\}\)
and the argument list of each existential \(z_j\) for \(\pi\) is \(z_1\)
followed by its argument list for \(\pi'\).

For part~(1), let \((f_{z_j})\) be a witness family for \(\pi\varphi\) at
\(s\).  Fix \(a\in M\) and put \(g^a_{z_j}(\bar u)=f_{z_j}(a,\bar u)\).
Every assignment of values to the universal variables of \(\pi'\), together
with \(z_1\mapsto a\), is an assignment of values to the universal variables
of \(\pi\), and the two families give each existential variable the same
value.  So \((g^a_{z_j})\) is a witness family for \(\pi'\varphi\) at
\(s[a/z_1]\).  By part~(1) for \(\pi'\),
\(M\entails\pi'\varphi\,[s[a/z_1]]\).  As \(a\) was arbitrary,
\(M\entails\forall z_1\pi'\varphi\,[s]\).

For part~(2), assume \(M\entails\forall z_1\pi'\varphi\,[s]\).  For each
\(a\in M\), \(M\entails\pi'\varphi\,[s[a/z_1]]\), so by part~(2) for
\(\pi'\) the set \(W_a\) of witness families for \(\pi'\varphi\) at
\(s[a/z_1]\) is nonempty.  By the Axiom of Choice there is a function
\(a\mapsto(g^a_{z_j})\in W_a\) on \(M\).  Define
\(f_{z_j}(a,\bar u)=g^a_{z_j}(\bar u)\).  Given values for the universal
variables of \(\pi\), let \(a\) be the value of \(z_1\).  The values the
family \((f_{z_j})\) assigns to the existential variables are those
assigned by \((g^a_{z_j})\) at the remaining universal values, and
\((g^a_{z_j})\in W_a\) makes \(\varphi\) true there.  So \((f_{z_j})\) is a
witness family for \(\pi\varphi\) at \(s\).

\emph{Case \(\mathsf Q_1=\exists\).}  No universal variable precedes
\(z_1\), so its argument list is empty and its witness is an element of
\(M\).  No pair of \(D_\pi\) has \(z_1\) as first entry, since \(z_1\) is
not universal.  Hence \(D_\pi=D_{\pi'}\), and every other existential
variable has the same argument list for \(\pi\) and for \(\pi'\).

For part~(1), let \((f_{z_j})\) be a witness family for \(\pi\varphi\) at
\(s\), and let \(a=f_{z_1}\).  The remaining functions form a witness family
for \(\pi'\varphi\) at \(s[a/z_1]\), because, for each choice of universal
values, they produce exactly the assignment that \((f_{z_j})\) produces
with \(z_1\mapsto a\).  By part~(1) for \(\pi'\),
\(M\entails\pi'\varphi\,[s[a/z_1]]\), so
\(M\entails\exists z_1\pi'\varphi\,[s]\).

For part~(2), choose \(a\in M\) with \(M\entails\pi'\varphi\,[s[a/z_1]]\).
Part~(2) for \(\pi'\) gives a witness family \((g_{z_j})\) for
\(\pi'\varphi\) at \(s[a/z_1]\).  Put \(f_{z_1}=a\) and \(f_{z_j}=g_{z_j}\)
for \(j\ge2\).  By the same identification of assignments, this is a
witness family for \(\pi\varphi\) at \(s\).  No choice is used in this case.

This completes the induction.  Only the satisfaction relation of
\(\varphi\) under assignments was used, so the argument applies to formulas
of \(\mathscr L_{\omega\omega}\) and of
\(\mathscr L_{\omega\omega}(Q^{\mathrm H})\) alike.
\end{proof}
\end{annexitem}

\begin{annexitem}{annex:p03-w03-p03-044-qh-vs-linear}{thm:p03-w03-p03-044-qh-vs-linear}{Removing a dependence strengthens the prefix}
\begin{proof}
(1) Let \(f_0,f_1:A\to A\) witness
\(\mathfrak A\entails Q^{\mathrm H}x_0y_0x_1y_1\,\varphi\,[s]\).  The
linear prefix \(\pi=\forall x_0\,\exists y_0\,\forall x_1\,\exists y_1\) has
argument list \((x_0)\) for \(y_0\) and \((x_0,x_1)\) for \(y_1\).  Put
\(g_{y_0}(a_0)=f_0(a_0)\) and \(g_{y_1}(a_0,a_1)=f_1(a_1)\).  For all
\(a_0,a_1\in A\), the values these give are \(f_0(a_0)\) and \(f_1(a_1)\),
and \(\varphi\) holds there by the choice of \(f_0,f_1\).  So
\((g_{y_0},g_{y_1})\) is a witness family for \(\pi\varphi\) at \(s\), and
part~(1) of \cref{thm:p03-w03-p03-044-linear-skolem}, which does not use
choice, gives \(\mathfrak A\entails\pi\varphi\,[s]\).

(2) Let \(A=\{0,1\}\) and \(\varphi\) be \(y_1=x_0\).  For the linear
sentence, given \(a_0\) choose \(y_0=0\); given then \(a_1\) choose
\(y_1=a_0\).  The matrix holds, so
\(\forall x_0\,\exists y_0\,\forall x_1\,\exists y_1\,(y_1=x_0)\) is true.
For the Henkin sentence, suppose \(f_0,f_1:A\to A\) were witnesses.  Taking
\(a_1=0\) and \(a_0=0\) gives \(f_1(0)=0\); taking \(a_1=0\) and \(a_0=1\)
gives \(f_1(0)=1\).  This is impossible, so
\(Q^{\mathrm H}x_0y_0x_1y_1\,(y_1=x_0)\) is false in \(A\).
\end{proof}
\end{annexitem}

\begin{annexitem}{annex:p03-w03-p03-044-qh-infinity}{thm:p03-w03-p03-044-qh-infinity}{The Henkin quantifier defines infinity}
\begin{proof}
Write \(\chi\) for the matrix
\(z\neq y_0\land(y_0=x_1\to y_1=x_0)\).

(1) By the clauses for \(\exists\) and for \(Q^{\mathrm H}\),
\(\mathfrak A\entails\varphi_\infty\) iff there are \(a\in A\) and
\(f_0,f_1:A\to A\) such that for all \(a_0,a_1\in A\):
\[
 (\ast)\qquad a\neq f_0(a_0)\quad\text{and}\quad
 \bigl(f_0(a_0)=a_1\ \Rightarrow\ f_1(a_1)=a_0\bigr).
\]
Assume \((\ast)\) for all \(a_0,a_1\).  The first conjunct, for every
\(a_0\), says \(a\notin\operatorname{rg}(f_0)\).  For \(b\in A\), apply
the second conjunct with \(a_0=b\) and \(a_1=f_0(b)\); its hypothesis holds,
so \(f_1(f_0(b))=b\).

Conversely, assume \(a\notin\operatorname{rg}(f_0)\) and \(f_1\circ
f_0=\mathrm{id}_A\).  For any \(a_0,a_1\), the first conjunct of \((\ast)\)
holds because \(f_0(a_0)\) lies in the range of \(f_0\) and \(a\) does not.
If \(f_0(a_0)=a_1\), then \(f_1(a_1)=f_1(f_0(a_0))=a_0\).  So \((\ast)\)
holds for all \(a_0,a_1\).

(2) Take \(a,f_0,f_1\) as in~(1).  If \(f_0(b)=f_0(b')\), then
\(b=f_1(f_0(b))=f_1(f_0(b'))=b'\), so \(f_0\) is one-to-one.  Suppose \(A\)
were finite with \(N\) elements.  The range of a one-to-one map on \(A\) has
\(N\) elements and is a subset of \(A\), so it equals \(A\).  Then \(a\) lies
in the range of \(f_0\), contrary to the choice of \(a\).  Hence \(A\) is
infinite.

(3) Let \(A\) be infinite and let \(c\) be a choice function on the
nonempty subsets of \(A\), supplied by the Axiom of Choice.  Define
\(e:\N\to A\) by recursion:
\(e(n)=c\bigl(A\setminus\{e(0),\ldots,e(n-1)\}\bigr)\).  The set removed is
finite and \(A\) is not, so the argument of \(c\) is nonempty at every
stage.  By construction \(e(n)\) differs from every earlier value, so \(e\)
is one-to-one.  Let \(R=\operatorname{rg}(e)\) and define
\[
 f_0(b)=\begin{cases} e(n+1), & b=e(n),\\ b, & b\notin R.\end{cases}
\]
This is well defined because \(e\) is one-to-one.  The map \(f_0\) sends
\(R\) into \(R\) one-to-one, since \(e(n+1)=e(m+1)\) gives \(n=m\), and it
is the identity on \(A\setminus R\).  The images of the two parts are
disjoint, so \(f_0\) is one-to-one on \(A\).  Put \(a=e(0)\).  It is not of
the form \(e(n+1)\), as \(e\) is one-to-one, and it does not lie in
\(f_0[A\setminus R]=A\setminus R\), since \(a\in R\).  So \(a\notin\operatorname{rg}(f_0)\).  Define \(f_1(d)\) to
be the unique \(b\) with \(f_0(b)=d\) when \(d\in\operatorname{rg}(f_0)\), and
\(f_1(d)=a\) otherwise.  Then \(f_1(f_0(b))=b\) for all \(b\in A\).  By~(1),
\(\mathfrak A\entails\varphi_\infty\).
\end{proof}
\end{annexitem}

\begin{annexitem}{annex:p03-w03-p03-044-qh-not-compact}{thm:p03-w03-p03-044-qh-not-compact}{\(\mathscr L_{\omega\omega}(Q^{\mathrm H})\) is not \(\aleph_0\)-compact}
\begin{proof}
For \(n\ge2\) let \(\sigma_n\) be the first-order sentence
\[
 \exists v_1\cdots\exists v_n\bigwedge_{1\le i<j\le n}v_i\neq v_j,
\]
which holds in a structure iff its universe has at least \(n\) elements.
Let
\[
 \Phi=\{\neg\varphi_\infty\}\cup\{\sigma_n:n\ge2\}.
\]
This is a countable set of \(\mathscr L_{\omega\omega}(Q^{\mathrm H})\)-sentences
of the empty vocabulary; \(\neg\varphi_\infty\) is a sentence because the
logic is closed under negation.

Let \(\Phi_0\subseteq\Phi\) be finite and let \(N\ge2\) exceed every \(n\)
with \(\sigma_n\in\Phi_0\).  The structure with universe
\(\{1,\ldots,N\}\) satisfies each \(\sigma_n\) with \(n\le N\).  It is
finite and nonempty, so by part~(2) of
\cref{thm:p03-w03-p03-044-qh-infinity} it does not satisfy
\(\varphi_\infty\); it satisfies \(\neg\varphi_\infty\).  So \(\Phi_0\) has a
model.

Suppose \(\mathfrak A\) were a model of \(\Phi\).  Its universe has at least
\(n\) elements for every \(n\ge2\), so it is nonempty and infinite.  By
part~(3) of \cref{thm:p03-w03-p03-044-qh-infinity}, which uses the Axiom of
Choice, \(\mathfrak A\entails\varphi_\infty\), contradicting
\(\mathfrak A\entails\neg\varphi_\infty\).  So \(\Phi\) has no model.

\(\Phi\) has power \(\aleph_0\), every finite subset has a model, and
\(\Phi\) has none.  By the definition of \(\kappa\)-compactness with
\(\kappa=\aleph_0\), the logic is not \(\aleph_0\)-compact.
\end{proof}
\end{annexitem}

\begin{annexitem}{annex:p03-w03-p03-044-qh-not-fo}{thm:p03-w03-p03-044-qh-not-fo}{No first-order equivalent of \(\varphi_\infty\)}
\begin{proof}
Suppose \(\theta\) were a first-order sentence of the empty vocabulary with
exactly the same nonempty models as \(\varphi_\infty\).  By
\cref{thm:p03-w03-p03-044-qh-infinity}, using the Axiom of Choice for
part~(3), a nonempty structure satisfies \(\varphi_\infty\) iff it is
infinite.  So a nonempty structure satisfies \(\neg\theta\) iff it is
finite, and the first-order theory \(\{\neg\theta\}\) has exactly the finite
nonempty structures of the empty signature as its models. For each
\(n\geq1\), let
\[
 \sigma_n=\exists x_1\cdots\exists x_n
             \bigwedge_{1\leq i<j\leq n}x_i\neq x_j
\]
assert that there are at least \(n\) elements (the conjunction for
\(n=1\) is true). Every finite subset of the first-order theory
\(T=\{\neg\theta\}\cup\{\sigma_n:n\geq1\}\) has a nonempty finite
model: take a finite set whose size is at least every mentioned \(n\).
First-order compactness (\cref{thm:fol-compactness}) gives a model of
\(T\). That model is infinite because it satisfies every \(\sigma_n\),
but finite because it satisfies \(\neg\theta\), a contradiction.
\end{proof}
\end{annexitem}

\begin{annexitem}{annex:p03-w03-p03-044-basic}{thm:p03-w03-p03-044-basic}{Basic properties of satisfaction}
\begin{proof}
Assignments are total, as in \cref{def:p03-w03-p03-044-henkin-syntax}.
Each part is proved by induction on \(\theta\), for all assignments
simultaneously, after the corresponding statement for terms, which is
proved by induction on terms.  For the step cases note that
\(s^{f_0f_1}_{a_0a_1}\) is obtained from \(s\) by four successive single
changes, so every identity between assignments used for \(\exists\) below
applies four times over in the \(Q^{\mathrm H}\) case.  The clauses for
\(\neg\) and \(\land\) refer to the immediate subformulas at the same
assignment, so in each part those cases follow at once from the induction
hypothesis.

(1) \emph{Terms.}  If \(s\) and \(s'\) agree at every variable of a term
\(t\), and every symbol of \(t\) is interpreted alike, then
\(t^{\mathfrak A}[s]=t^{\mathfrak B}[s']\), by induction on \(t\).  In the
atomic case every variable of the formula is free, so this applies to its
terms.  Atomic formulas then have
the same truth value, since their relation symbols, and equality, are
interpreted alike over the common universe.  \(\exists x\,\theta_1\): for
each \(a\), the assignments \(s[a/x]\) and \(s'[a/x]\) agree on the free
variables of \(\theta_1\), which are those of \(\exists x\,\theta_1\)
together possibly with \(x\); the induction hypothesis applies for each
\(a\) in the common universe.  \(Q^{\mathrm H}x_0y_0x_1y_1\,\theta_1\): for
the same pair \(f_0,f_1\) on the common universe and the same
\(a_0,a_1\), the assignments \(s^{f_0f_1}_{a_0a_1}\) and
\(s'^{f_0f_1}_{a_0a_1}\) agree on the free variables of \(\theta_1\), so the
same pairs are witnesses on both sides.

(2) \emph{Terms.}  \(t^{\mathfrak B}[\pi\circ s]=\pi(t^{\mathfrak A}[s])\),
because \(\pi\) maps \(c^{\mathfrak A}\) to \(c^{\mathfrak B}\) and commutes
with each \(f\).  Atomic formulas: \(\pi\) is one-to-one and preserves and
reflects each relation.  \(\exists x\,\theta_1\): since
\((\pi\circ s)[\pi(a)/x]=\pi\circ(s[a/x])\) and \(\pi\) is onto, a witness
\(a\) on the left gives the witness \(\pi(a)\) on the right, and a witness
\(b\) on the right gives \(\pi^{-1}(b)\) on the left.
\(Q^{\mathrm H}x_0y_0x_1y_1\,\theta_1\): let \(f_0,f_1\) be witnesses on the
left and put \(g_i=\pi\circ f_i\circ\pi^{-1}:B\to B\).  For
\(b_0,b_1\in B\) put \(a_i=\pi^{-1}(b_i)\); then \(g_i(b_i)=\pi(f_i(a_i))\)
and
\[
 (\pi\circ s)^{g_0g_1}_{b_0b_1}=\pi\circ\bigl(s^{f_0f_1}_{a_0a_1}\bigr),
\]
so the induction hypothesis turns
\(\mathfrak A\entails\theta_1[s^{f_0f_1}_{a_0a_1}]\) into
\(\mathfrak B\entails\theta_1[(\pi\circ s)^{g_0g_1}_{b_0b_1}]\).  The
converse direction is the same argument for \(\pi^{-1}:\mathfrak B\cong
\mathfrak A\), applied to the assignment \(\pi\circ s\).

(3) \(\mathfrak A^\rho\) has the universe of \(\mathfrak A\) and
\(\rho(S)^{\mathfrak A^\rho}=S^{\mathfrak A}\), so
\((t^\rho)^{\mathfrak A^\rho}[s]=t^{\mathfrak A}[s]\) for terms and atomic
formulas keep their truth values.  The quantifier clauses range over the
same universe and the same functions on it, and \(\theta^\rho\) has the same
variables as \(\theta\), so each step is immediate from the induction
hypothesis.

(4) \emph{Terms.}  \((\iota t)^{\mathfrak A}[s]=t^{\mathfrak A}[s\circ\iota]\),
since a variable \(v\) is sent to \(\iota(v)\), whose value under \(s\) is
\((s\circ\iota)(v)\).  Atomic formulas follow.  \(\exists x\,\theta_1\):
\(\iota(\exists x\,\theta_1)=\exists\iota(x)\,\iota\theta_1\), and for each
\(a\in A\)
\[
 (s[a/\iota(x)])\circ\iota=(s\circ\iota)[a/x],
\]
because both send \(x\) to \(a\), and for \(v\neq x\) injectivity gives
\(\iota(v)\neq\iota(x)\), so both send \(v\) to \(s(\iota(v))\).  With the
induction hypothesis, for each \(a\),
\[
 \mathfrak A\entails\iota\theta_1[s[a/\iota(x)]]\iff
 \mathfrak A\entails\theta_1[(s\circ\iota)[a/x]].
\]
\(Q^{\mathrm H}x_0y_0x_1y_1\,\theta_1\): \(\iota\) sends it to
\(Q^{\mathrm H}\iota(x_0)\iota(y_0)\iota(x_1)\iota(y_1)\,\iota\theta_1\),
whose four variables are distinct because \(\iota\) is injective; the same
identity, applied four times, gives
\(s^{f_0f_1}_{a_0a_1}\circ\iota=(s\circ\iota)^{f_0f_1}_{a_0a_1}\), where on
the left the changes are made at \(\iota(x_0),\ldots,\iota(y_1)\).  So the
same pairs \(f_0,f_1\) are witnesses on both sides.

(5) No \(v_i\) occurs in any subformula of \(\theta\), so the hypothesis is
inherited.  \emph{Terms.}  \((t(\bar v/\bar c))^{\mathfrak A\restriction
\tau_0}[s]=t^{\mathfrak A}[s]\): the occurrence of \(c_i\) in \(t\) becomes
\(v_i\), whose value \(s(v_i)\) is \(c_i^{\mathfrak A}\); a variable of \(t\)
is not a \(v_i\) and is unchanged; the other symbols lie in \(\tau_0\) and
are interpreted alike.  Atomic formulas follow, their relation symbols lying
in \(\tau_0\).  \(\exists x\,\theta_1\): the variable \(x\) occurs in
\(\theta\), so it is no \(v_i\), and \(s[a/x](v_i)=s(v_i)=c_i^{\mathfrak A}\)
for every \(a\); the induction hypothesis applies to \(s[a/x]\).
\(Q^{\mathrm H}x_0y_0x_1y_1\,\theta_1\): likewise none of the four bound
variables is a \(v_i\), so \(s^{f_0f_1}_{a_0a_1}(v_i)=c_i^{\mathfrak A}\),
the two structures have the same universe and the same functions on it, and
the same pairs are witnesses on both sides.
\end{proof}
\end{annexitem}

\begin{annexitem}{annex:p03-w03-p03-044-unnesting}{thm:p03-w03-p03-044-unnesting}{Every formula has a flat equivalent}
\begin{proof}
Replace \(W\) by \(W\cup\operatorname{var}(\theta)\), where
\(\operatorname{var}(\theta)\) is the finite set of variables occurring in
\(\theta\); new variables are then chosen outside the enlarged \(W\), which
gives the stated condition on variables.

\emph{Atomic formulas.}  For atomic \(\alpha\) let \(k(\alpha)\) be the
number of occurrences of constants and function symbols in \(\alpha\).  We
define \(\alpha^{\mathrm u}\), for every finite \(W\supseteq\operatorname
{var}(\alpha)\), by strong induction on \(k(\alpha)\).  If \(\alpha\) is
flat, put \(\alpha^{\mathrm u}=\alpha\).  Otherwise \(\alpha\) contains a
term that is not a variable.  Every such term has a subterm occurrence
\(t=hv_1\cdots v_n\) (\(n\ge0\); \(h\) a function symbol, or a constant
when \(n=0\)) whose arguments are variables: descend into an argument that
is not a variable as long as there is one.  Fix one such occurrence, let
\(z\) be a variable outside \(W\), and let \(\alpha[z/t]\) be the result of
replacing this occurrence of \(t\) by \(z\).  Then
\(k(\alpha[z/t])=k(\alpha)-1\), and we put
\[
 \alpha^{\mathrm u}=\exists z\,\bigl(hv_1\cdots v_n=z\land
 (\alpha[z/t])^{\mathrm u}\bigr),
\]
the inner flat formula being formed for the set \(W\cup\{z\}\).  The
conjunct \(hv_1\cdots v_n=z\) is flat, so \(\alpha^{\mathrm u}\) is flat.
Its free variables are those of \(\alpha\).  By the induction hypothesis
\((\alpha[z/t])^{\mathrm u}\) has the free variables of \(\alpha[z/t]\).
These are the variables of \(\alpha\) other than those that occurred only
in the replaced occurrence of \(t\), together with \(z\).  The conjunct
\(hv_1\cdots v_n=z\) restores \(v_1,\ldots,v_n\), and \(z\) is bound.

Equivalence.  Fix \(\mathfrak A\) and \(s\), and let
\(b=t^{\mathfrak A}[s]\).  Because \(z\notin\operatorname{var}(\alpha)\), an
induction on terms shows that every term of \(\alpha[z/t]\) has at
\(s[b/z]\) the value of the corresponding term of \(\alpha\) at \(s\): the
new occurrence of \(z\) has value \(b\), and every other variable keeps its
value.  Hence \(\mathfrak A\entails\alpha[z/t]\,[s[b/z]]\) iff
\(\mathfrak A\entails\alpha\,[s]\).  Also, for \(b'\in A\),
\(\mathfrak A\entails hv_1\cdots v_n=z\,[s[b'/z]]\) iff \(b'=b\), since
\(z\) is none of the \(v_i\).  So
\(\mathfrak A\entails\exists z\,(hv_1\cdots v_n=z\land\alpha[z/t])\,[s]\)
iff \(\mathfrak A\entails\alpha\,[s]\), and by the induction hypothesis
\((\alpha[z/t])^{\mathrm u}\) may replace \(\alpha[z/t]\) inside, giving
\(\mathfrak A\entails\alpha^{\mathrm u}[s]\) iff
\(\mathfrak A\entails\alpha[s]\).

\emph{Compound formulas.}  Put \((\neg\theta_1)^{\mathrm u}=\neg
\theta_1^{\mathrm u}\), \((\theta_0\land\theta_1)^{\mathrm u}=
\theta_0^{\mathrm u}\land\theta_1^{\mathrm u}\), \((\exists x\,
\theta_1)^{\mathrm u}=\exists x\,\theta_1^{\mathrm u}\) and
\((Q^{\mathrm H}x_0y_0x_1y_1\,\theta_1)^{\mathrm u}=Q^{\mathrm
H}x_0y_0x_1y_1\,\theta_1^{\mathrm u}\), all formed for the same \(W\).
Free variables are preserved because they are preserved in the immediate
subformulas.  Each satisfaction clause refers to the immediate subformulas
only through their satisfaction at assignments over the same structure, and
by the induction hypothesis these agree at every assignment; so the
equivalence passes to the compound formula.  This holds for the
\(Q^{\mathrm H}\) clause in particular: a pair \(f_0,f_1\) witnesses
\(\theta_1\) at \(s\) iff it witnesses \(\theta_1^{\mathrm u}\) at \(s\).
\end{proof}
\end{annexitem}

\begin{annexitem}{annex:p03-w03-p03-044-regular}{thm:p03-w03-p03-044-regular}{Regularity of the Henkin logic}
\begin{proof}
Write \(\mathscr L\) for \(\mathscr L_{\omega\omega}(Q^{\mathrm H})\) and
\(\operatorname{var}(\theta)\) for the set of variables occurring in
\(\theta\).  By \cref{thm:p03-w03-p03-044-basic}(1), the truth of a sentence
does not depend on the assignment.  If \(\psi\) is a sentence and \(\iota\)
is injective, \(\iota\psi\) is again a sentence and, by
\cref{thm:p03-w03-p03-044-basic}(4) and~(1),
\(\mathfrak A\entails\iota\psi\) iff \(\mathfrak A\entails\psi\).  Since
there are infinitely many variables, for every sentence \(\psi\) and every
finite set \(V\) of variables there is such an \(\iota\) for which
\(\operatorname{var}(\iota\psi)\) is disjoint from \(V\).

(1) (i): a sentence all of whose symbols lie in \(\tau\) has all its
symbols in any \(\sigma\supseteq\tau\).  (ii): \(\mathfrak A\entails
\varphi\) is defined only for \(\varphi\in\mathscr L[\tau_{\mathfrak A}]\).
(iii): by \cref{thm:p03-w03-p03-044-basic}(2) for a sentence.  (iv): if
\(\varphi\in\mathscr L[\tau]\) and \(\tau\subseteq\tau_{\mathfrak A}\), then
\(\mathfrak A\) and \(\mathfrak A\restriction\tau\) have the same universe
and interpret the symbols of \(\varphi\) alike, so
\cref{thm:p03-w03-p03-044-basic}(1) applies.  (v): \(\varphi^\rho\) is a
sentence of \(\sigma\), and \cref{thm:p03-w03-p03-044-basic}(3) applies.

(2) \emph{Atom.}  An atomic first-order sentence \(\varphi\) of \(\tau\) lies
in \(\mathscr L[\tau]\), and its clause in
\cref{def:p03-w03-p03-044-henkin-syntax} is the first-order one; take
\(\psi=\varphi\).  \emph{Negation and conjunction.}  Take \(\neg\varphi\) and
\(\varphi_0\land\varphi_1\); their clauses give the complement and the
intersection of the model classes.  \emph{Particularization.}  Let
\(c\in\tau\), \(\varphi\in\mathscr L[\tau]\), and \(x\notin
\operatorname{var}(\varphi)\).  Then \(\psi=\exists x\,\varphi(x/c)\) is a
sentence of \(\tau\setminus\{c\}\).  For a \((\tau\setminus\{c\})\)-structure
\(\mathfrak A\) and any \(s\), \(\mathfrak A\entails\psi\) iff
\(\mathfrak A\entails\varphi(x/c)\,[s[a/x]]\) for some \(a\in A\).  For each
\(a\), apply \cref{thm:p03-w03-p03-044-basic}(5) to the \(\tau\)-structure
\((\mathfrak A,a)\), in which \(c\) denotes \(a\), and the assignment
\(s[a/x]\), which sends \(x\) to \(a\): it gives
\(\mathfrak A\entails\varphi(x/c)\,[s[a/x]]\) iff
\((\mathfrak A,a)\entails\varphi\).  So \(\mathfrak A\entails\psi\) iff
\((\mathfrak A,a)\entails\varphi\) for some \(a\in A\).

(3) \emph{Relativization.}  Let \(c\notin\tau\cup\sigma\),
\(\chi\in\mathscr L[\sigma\cup\{c\}]\) and \(\varphi\in\mathscr L[\tau]\),
and let \(V=\operatorname{var}(\varphi)\).  Choose \(\iota\) with
\(\operatorname{var}(\iota\chi)\cap V=\varnothing\) and put
\(\chi_v=(\iota\chi)(v/c)\) for \(v\in V\); its only free variable is
\(v\).  For a \((\tau\cup\sigma)\)-structure \(\mathfrak B\) and an
assignment \(s\), apply \cref{thm:p03-w03-p03-044-basic}(5) to
\((\mathfrak B,s(v))\) and then (4), (1):
\[
 \mathfrak B\entails\chi_v\,[s]\iff(\mathfrak B,s(v))\entails\iota\chi
 \iff(\mathfrak B,s(v))\entails\chi\iff s(v)\in\chi^{\mathfrak B}.
\]
Define \(\theta^U\) for formulas \(\theta\) of \(\tau\) with
\(\operatorname{var}(\theta)\subseteq V\): atomic formulas are unchanged,
\((\neg\theta_1)^U=\neg\theta_1^U\),
\((\theta_0\land\theta_1)^U=\theta_0^U\land\theta_1^U\),
\((\exists x\,\theta_1)^U=\exists x\,(\chi_x\land\theta_1^U)\), and
\[
 (Q^{\mathrm H}x_0y_0x_1y_1\,\theta_1)^U=
 Q^{\mathrm H}x_0y_0x_1y_1\bigl((\chi_{x_0}\land\chi_{x_1})\to
 (\chi_{y_0}\land\chi_{y_1}\land\theta_1^U)\bigr).
\]
Each \(\chi_v\) is inserted inside the scope of the quantifier binding
\(v\), and its bound variables lie outside \(V\), so the free variables of
\(\theta^U\) are among those of \(\theta\).  Put \(\psi=\varphi^U\); it is a
sentence of \(\tau\cup\sigma\).

Let \(C=\chi^{\mathfrak B}\) be \(\tau\)-closed in \(\mathfrak B\), that is,
nonempty, containing every constant of \(\tau\) and closed under every
function of \(\tau\) (\cref{def:p03-w03-p03-047-vocabularies}), and let
\(\mathfrak C=(\mathfrak B\restriction\tau)\,|\,C\).  We show, by induction
on \(\theta\), that \(\mathfrak C\entails\theta\,[s]\) iff
\(\mathfrak B\entails\theta^U\,[s]\) for every assignment \(s\) into \(C\).
\emph{Terms}: for \(s\) into \(C\), \(t^{\mathfrak B}[s]=t^{\mathfrak C}[s]\in
C\), by induction on \(t\), because \(C\) contains the constants and is
closed under the functions, which \(\mathfrak C\) interprets by
restriction.  Atomic formulas follow, as \(\mathfrak C\) carries the
restricted relations.  \(\neg\), \(\land\): immediate.
\(\exists x\,\theta_1\): \(\mathfrak B\entails(\exists x\,\theta_1)^U[s]\)
iff some \(b\in B\) has \(\mathfrak B\entails\chi_x[s[b/x]]\), that is
\(b\in C\), and \(\mathfrak B\entails\theta_1^U[s[b/x]]\); as \(s[b/x]\)
is into \(C\), the induction hypothesis turns this into
\(\mathfrak C\entails\theta_1[s[b/x]]\) for some \(b\in C\), that is
\(\mathfrak C\entails\exists x\,\theta_1[s]\).

\(Q^{\mathrm H}x_0y_0x_1y_1\,\theta_1\), from \(\mathfrak C\) to
\(\mathfrak B\): let \(f_0,f_1:C\to C\) be witnesses in \(\mathfrak C\).
Fix \(e\in C\) and put \(g_i(b)=f_i(b)\) for \(b\in C\) and \(g_i(b)=e\)
otherwise.  Let \(b_0,b_1\in B\).  If \(b_0\notin C\) or \(b_1\notin C\),
then \(\chi_{x_0}\land\chi_{x_1}\) fails at \(s^{g_0g_1}_{b_0b_1}\) and the
relativized matrix holds there.  If \(b_0,b_1\in C\), then
\(g_i(b_i)=f_i(b_i)\in C\), so \(\chi_{y_0}\) and \(\chi_{y_1}\) hold, the
assignment \(s^{g_0g_1}_{b_0b_1}=s^{f_0f_1}_{b_0b_1}\) is into \(C\), and the
induction hypothesis turns \(\mathfrak C\entails\theta_1[s^{f_0f_1}_{b_0b_1}]\)
into \(\mathfrak B\entails\theta_1^U[s^{g_0g_1}_{b_0b_1}]\).  So
\(g_0,g_1\) are witnesses in \(\mathfrak B\).

From \(\mathfrak B\) to \(\mathfrak C\): let \(g_0,g_1:B\to B\) be witnesses
for the relativized matrix.  Let \(b_0\in C\); since \(C\neq\varnothing\),
pick \(b_1\in C\).  At \(s^{g_0g_1}_{b_0b_1}\) the hypothesis
\(\chi_{x_0}\land\chi_{x_1}\) holds, so \(\chi_{y_0}\) holds, that is
\(g_0(b_0)\in C\).  In the same way \(g_1(b_1)\in C\) for every \(b_1\in
C\).  Let \(f_i\) be the restriction of \(g_i\) to \(C\), a map \(C\to C\).
For \(a_0,a_1\in C\) the hypothesis holds at
\(s^{g_0g_1}_{a_0a_1}=s^{f_0f_1}_{a_0a_1}\), which is into \(C\), so
\(\mathfrak B\entails\theta_1^U\) there, and by the induction hypothesis
\(\mathfrak C\entails\theta_1[s^{f_0f_1}_{a_0a_1}]\).  So \(f_0,f_1\) are
witnesses in \(\mathfrak C\).  Neither direction uses choice: \(e\),
\(b_1\) are single elements fixed once.

Applying the claim to the sentence \(\varphi\) and any assignment into the
nonempty set \(C\) gives \(\mathfrak B\entails\psi\) iff
\(\mathfrak C\entails\varphi\), which is the relativization property.

(4) \emph{Substitution.}  Let \(\varphi^{\mathrm u}\) be the flat
equivalent of \(\varphi\) from \cref{thm:p03-w03-p03-044-unnesting} (with
\(W=\varnothing\)); it is a sentence with the same symbols as \(\varphi\),
since unnesting introduces no symbol.  Let \(V=\operatorname{var}
(\varphi^{\mathrm u})\).  For each \(j\le k\) choose an injective
\(\iota_j\) with \(\psi'_j=\iota_j\psi_{S_j}\) having no variable in \(V\),
and for a tuple \(\bar v\) of variables, of the length of \(\mathbf
c_{S_j}\) and none occurring in \(\psi'_j\), put
\(\psi_j^{\bar v}=\psi'_j(\bar v/\mathbf c_{S_j})\), a formula of \(\tau\)
whose free variables are among \(\bar v\).  For a \(\tau\)-structure
\(\mathfrak A\) and assignment \(s\), \cref{thm:p03-w03-p03-044-basic}(5)
applied to the expansion of \(\mathfrak A\) interpreting \(\mathbf c_{S_j}\)
by \(s(\bar v)\), followed by~(4) and~(1), gives
\[
 (\dagger)\qquad
 \mathfrak A\entails\psi_j^{\bar v}\,[s]\iff
 (\mathfrak A,s(\bar v))\entails\psi_{S_j}.
\]
Repeated variables in \(\bar v\) are allowed in~(5), the hypothesis there
being that \(s(v_i)\) is the value of the \(i\)-th constant.

Define \(T\theta\) for flat formulas \(\theta\) of \(\tau'\) with
\(\operatorname{var}(\theta)\subseteq V\):
\(T(S_j\bar v)=\psi_j^{\bar v}\) for a relation symbol \(S_j\);
\(T(S_j\bar v=w)=\psi_j^{\bar vw}\) for a function symbol \(S_j\);
\(T(S_j=w)=\psi_j^{w}\) for a constant \(S_j\); every other atomic formula
is unchanged; and \(T\) commutes with \(\neg\), \(\land\), \(\exists x\) and
\(Q^{\mathrm H}x_0y_0x_1y_1\).  An atomic formula left unchanged contains
only symbols of \(\tau\): a symbol of \(\tau'\setminus\tau\) occurring in
\(\theta\) is some \(S_j\), and in a flat formula it occurs only in one of
the three forms treated.  So \(T\theta\) is a formula of \(\tau\), and its
free variables are among those of \(\theta\).

For each function symbol \(S_j\) of arity \(n\) let \(\mathrm{Fun}_j\) be
\(\forall\bar x\,\exists y\,\forall y'(\psi_j^{\bar xy'}\leftrightarrow
y'=y)\), and for each constant \(S_j\) let \(\mathrm{Fun}_j\) be
\(\exists y\,\forall y'(\psi_j^{y'}\leftrightarrow y'=y)\), with
\(\bar x,y,y'\) distinct variables not occurring in \(\psi'_j\).  By
\((\dagger)\), \(\mathfrak A\entails\mathrm{Fun}_j\) iff the set defined by
\(\psi_{S_j}\) is the graph of a function \(A^n\to A\), or a singleton
\(\{a\}\), whose single element is then taken as the value of the constant.
Let \(\varphi^\ast\) be the conjunction of all \(\mathrm{Fun}_j\) and
\(T\varphi^{\mathrm u}\); it lies in \(\mathscr L[\tau]\).

Let \(\mathfrak A\) be a \(\tau\)-structure.  If some \(\mathrm{Fun}_j\)
fails, no expansion of the required kind exists and
\(\mathfrak A\not\entails\varphi^\ast\).  Otherwise let \(\mathfrak A_1\) be
the unique expansion of \(\mathfrak A\) to \(\tau\cup\{S_1,\ldots,S_k\}\)
interpreting each \(S_j\) as required.  We show
\(\mathfrak A_1\entails\theta\,[s]\) iff \(\mathfrak A\entails
T\theta\,[s]\) for flat \(\theta\) with variables in \(V\) and symbols in
\(\tau\cup\{S_1,\ldots,S_k\}\), by induction on \(\theta\).  For
\(S_j\bar v\): \(\mathfrak A_1\entails S_j\bar v[s]\) iff \(s(\bar v)\in
S_j^{\mathfrak A_1}\) iff \((\mathfrak A,s(\bar v))\entails\psi_{S_j}\) iff
\(\mathfrak A\entails\psi_j^{\bar v}[s]\) by \((\dagger)\).  For
\(S_j\bar v=w\): it holds iff \((s(\bar v),s(w))\) lies in the graph of
\(S_j^{\mathfrak A_1}\), which is the set defined by \(\psi_{S_j}\); by
\((\dagger)\) this is \(\mathfrak A\entails\psi_j^{\bar vw}[s]\).  The
constant case is the same with \(n=0\).  Unchanged atomic formulas have
symbols in \(\tau\), on which \(\mathfrak A_1\) and \(\mathfrak A\) agree,
so \cref{thm:p03-w03-p03-044-basic}(1) applies.  The step cases hold
because \(\mathfrak A_1\) and \(\mathfrak A\) have the same universe, so
the clauses for \(\exists\) and \(Q^{\mathrm H}\) range over the same
elements and the same pairs of functions, at the same assignments.

Hence \(\mathfrak A\entails\varphi^\ast\) iff \(\mathfrak A_1\entails
\varphi^{\mathrm u}\) iff \(\mathfrak A_1\entails\varphi\), by
\cref{thm:p03-w03-p03-044-unnesting}.  Every expansion \(\mathfrak A'\) of
\(\mathfrak A_1\) to \(\tau'\) satisfies \(\varphi\) iff \(\mathfrak A_1\)
does, by \cref{thm:p03-w03-p03-044-basic}(1), and such an expansion exists
because \(A\neq\varnothing\).  The expansions of \(\mathfrak A\) to \(\tau'\)
of the required kind are exactly these.  So \(\mathfrak A\entails
\varphi^\ast\) iff \(\mathfrak A\) has an expansion of the required kind
satisfying \(\varphi\).

If only finitely many symbols of \(\tau'\setminus\tau\) are function
symbols or constants, add to \(\varphi^\ast\) the finitely many sentences
\(\mathrm{Fun}_S\), formed in the same way, for the function symbols and
constants \(S\in\tau'\setminus\tau\) that do not occur in \(\varphi\).
Then \(\mathfrak A\) satisfies the enlarged sentence iff every such
\(\psi_S\) defines a graph and \(\mathfrak A\entails\varphi^\ast\); and in
that case the expansion of \(\mathfrak A_1\) that interprets every
remaining symbol as \(\psi_S\) requires satisfies \(\varphi\) iff
\(\mathfrak A_1\) does, by \cref{thm:p03-w03-p03-044-basic}(1).  Relation
symbols that do not occur need no conjunct, since every set is a relation;
there may be infinitely many of them, and each is interpreted as the set
that its \(\psi_S\) defines.

Parts (1)--(4) are the conditions of
\cref{def:p03-w03-p03-047-ebbinghaus-logic,def:p03-w03-p03-047-closure,def:p03-w03-p03-047-relativization}
on one-sorted vocabularies.  Choice entered nowhere: every function used is
defined explicitly from given functions and finitely many fixed elements.
\end{proof}
\end{annexitem}

\begin{annexitem}{annex:p03-w03-p03-044-substitution-limit}{thm:p03-w03-p03-044-substitution-limit}{Substitution of infinitely many constants}
\begin{proof}
Write \(\mathscr L\) for \(\mathscr L_{\omega\omega}(Q^{\mathrm H})\) and
\(N_2=\{n\in\mathbb N:n\ge2\}\).  Every structure has a nonempty universe.

\emph{Step 1: the sentences \(\varepsilon_n\).}  Let \(\mathfrak A\) be a
structure and \(n\ge2\).  If \(\mathfrak A\entails\varepsilon_n\), there are
\(a_1,\ldots,a_n\in A\), pairwise distinct, such that every \(a\in A\)
equals some \(a_i\); so \(A=\{a_1,\ldots,a_n\}\) has exactly \(n\) elements.
Conversely, if \(A=\{a_1,\ldots,a_n\}\) with the \(a_i\) pairwise distinct,
these values witness the existential quantifiers of \(\varepsilon_n\).  So
\(\mathfrak A\entails\varepsilon_n\) iff \(|A|=n\).

\emph{Step 2: the required class, for an arbitrary
\(X\subseteq N_2\).}  Fix \(X\), \(\tau'\), \(\varphi\) and the
\(\psi_{c_n}\) as in the statement, and let \(\mathfrak A\) be a
\(\tau\)-structure.  For \(n\in X\) put
\(D_n=\{a\in A:(\mathfrak A,a)\entails\psi_{c_n}\}\), where \(d\) is
interpreted by \(a\).  Since \((\mathfrak A,a)\entails d=e\) iff
\(a=e^{\mathfrak A}\), and by Step 1 \((\mathfrak A,a)\entails\neg
\varepsilon_n\) iff \(|A|\neq n\), we get \(D_n=\{e^{\mathfrak A}\}\) if
\(|A|\neq n\) and \(D_n=\varnothing\) if \(|A|=n\).

If \(|A|\notin X\) (this includes every infinite \(\mathfrak A\)), let
\(\mathfrak A'\) be the expansion of \(\mathfrak A\) to \(\tau'\) with
\(c_n^{\mathfrak A'}=e^{\mathfrak A}\) for every \(n\in X\).  For each
\(n\in X\) we have \(n\neq|A|\), so \(\{c_n^{\mathfrak A'}\}=
\{e^{\mathfrak A}\}=D_n\); and \(\mathfrak A'\entails e=e\).  So
\(\mathfrak A\) has an expansion of the required kind.  If \(|A|=m\in X\),
then \(D_m=\varnothing\), while \(\{c_m^{\mathfrak A'}\}\) has one element
for every expansion \(\mathfrak A'\); so no expansion of the required kind
exists.  Hence the \(\tau\)-structures that have an expansion of the
required kind form the class
\[
 K_X=\{\mathfrak A\in\operatorname{Str}[\tau]:|A|\notin X\}.
\]
A sentence \(\varphi^\ast\in\mathscr L[\tau]\) has the property of the
statement exactly when its \(\tau\)-models are the members of \(K_X\).

\emph{Step 3: finite sizes.}  For \(m\ge2\) let \(\mathfrak M_m\) be the
\(\tau\)-structure with universe \(\{0,\ldots,m-1\}\) and
\(e^{\mathfrak M_m}=0\).  For \(\chi\in\mathscr L[\tau]\) put
\(S_\chi=\{m\in N_2:\mathfrak M_m\entails\chi\}\).  If \(\varphi^\ast\) has
the property of the statement for \(X\), then by Step 2, for \(m\in N_2\),
\(\mathfrak M_m\entails\varphi^\ast\) iff \(\mathfrak M_m\in K_X\) iff
\(m\notin X\); that is, \(S_{\varphi^\ast}=N_2\setminus X\).

\emph{Step 4: counting.}  Each sentence of \(\mathscr L[\tau]\) is a finite
string over the alphabet consisting of the variables, the constant \(e\),
the symbols \(=\), \(\neg\), \(\land\), \(\exists\) and \(Q^{\mathrm H}\), and
parentheses (\cref{def:p03-w03-p03-044-henkin-syntax}).  This alphabet is
countable because the set of variables is, and the set of finite strings
over a countable alphabet is the countable union over \(k\in\mathbb N\) of
the countable sets of strings of length \(k\).  So \(\mathscr L[\tau]\) is
countable, and so is the set \(\{S_\chi:\chi\in\mathscr L[\tau]\}\).  The
map \(X\mapsto N_2\setminus X\) is a bijection of the power set of \(N_2\)
onto itself, and by Cantor's theorem this power set is uncountable.  Hence
there is \(X\subseteq N_2\) such that \(N_2\setminus X\neq S_\chi\) for
every \(\chi\in\mathscr L[\tau]\).  By Step 3 no
\(\varphi^\ast\in\mathscr L[\tau]\) has the property of the statement for
this \(X\).

\emph{First-order logic.}  The first-order sentences of vocabulary \(\tau\)
are sentences of \(\mathscr L[\tau]\) without \(Q^{\mathrm H}\), and the
clauses of \cref{def:p03-w03-p03-044-henkin-syntax} for them are the
first-order clauses.  A first-order \(\varphi^\ast\) with the property would
thus be a sentence of \(\mathscr L[\tau]\) with the property, which Step 4
excludes.  No use of the Axiom of Choice is made.
\end{proof}
\end{annexitem}

\begin{annexitem}{annex:p03-w03-p03-044-flatness}{thm:p03-w03-p03-044-flatness}{First-order formulas are flat on teams}
\begin{proof}
By induction on \(\psi\), for all teams \(X\) simultaneously, using the
clauses of \cref{def:p03-team-semantics}.

\emph{First-order literal.}  The team clause is the right side.

\emph{\(\bottom\).}  \(M,X\entails\bottom\) iff \(X=\varnothing\).  No
assignment satisfies \(\bottom\), so \((\forall s\in X)\,M,s\entails\bottom\)
also holds iff \(X=\varnothing\).

\emph{\(\psi_1\land\psi_2\).}  Both sides split into the two conjuncts, and
the induction hypothesis applies to each.

\emph{\(\psi_1\lor\psi_2\).}  Suppose \(X=Y\cup Z\) with
\(M,Y\entails\psi_1\) and \(M,Z\entails\psi_2\).  By the induction
hypothesis every \(s\in Y\) satisfies \(\psi_1\) and every \(s\in Z\)
satisfies \(\psi_2\), so every \(s\in X\) satisfies \(\psi_1\lor\psi_2\).
Conversely, if every \(s\in X\) satisfies \(\psi_1\lor\psi_2\), put
\(Y=\{s\in X:M,s\entails\psi_1\}\) and \(Z=\{s\in X:M,s\entails\psi_2\}\).
Then \(X=Y\cup Z\), and the induction hypothesis gives
\(M,Y\entails\psi_1\) and \(M,Z\entails\psi_2\).

\emph{\(\exists x\,\psi_1\).}  Suppose \(M,X[F/x]\entails\psi_1\) for some
\(F:X\to\mathcal P(M)\setminus\{\varnothing\}\).  For \(s\in X\) pick
\(a\in F(s)\); then \(s[a/x]\in X[F/x]\), so by the induction hypothesis
\(M,s[a/x]\entails\psi_1\) and \(M,s\entails\exists x\,\psi_1\).  No choice
is needed, because one element is picked for one fixed \(s\).  Conversely,
if every \(s\in X\) satisfies \(\exists x\,\psi_1\), put
\(F(s)=\{a\in M:M,s[a/x]\entails\psi_1\}\), which is nonempty.  Every member
of \(X[F/x]\) satisfies \(\psi_1\), so the induction hypothesis gives
\(M,X[F/x]\entails\psi_1\).

\emph{\(\forall x\,\psi_1\).}  By the induction hypothesis,
\(M,X[M/x]\entails\psi_1\) iff \(M,s[a/x]\entails\psi_1\) for all \(s\in X\)
and \(a\in M\), that is, iff every \(s\in X\) satisfies \(\forall x\,\psi_1\).
\end{proof}
\end{annexitem}

\begin{annexitem}{annex:p03-w03-p03-044-qh-team}{thm:p03-w03-p03-044-qh-team}{The Henkin quantifier as a dependence-logic sentence}
\begin{proof}
For \(a,b,c,d\in A\) write
\(s_{a}=s[a/x_0]\), \(s_{ab}=s_a[b/y_0]\), \(s_{abc}=s_{ab}[c/x_1]\) and
\(s_{abcd}=s_{abc}[d/y_1]\).  Unwinding the team clauses of
\cref{def:p03-team-semantics}, the right side holds iff there are
\(F_0\) on \(X_1=\{s_a:a\in A\}\) and \(F_1\) on
\[
 X_3=\{s_{abc}:a\in A,\ b\in F_0(s_a),\ c\in A\},
\]
both with nonempty values, such that the team
\[
 X_4=\{s_{abcd}:s_{abc}\in X_3,\ d\in F_1(s_{abc})\}
\]
satisfies \(=(x_1,y_1)\land\psi\).  By the conjunction clause and
\cref{thm:p03-w03-p03-044-flatness}, this means that any two members of
\(X_4\) with the same value at \(x_1\) have the same value at \(y_1\), and
that every member of \(X_4\) satisfies \(\psi\).

\emph{Left to right.}  Let \(f_0,f_1\) witness the Henkin side.  Put
\(F_0(s_a)=\{f_0(a)\}\) and \(F_1(s_{abc})=\{f_1(c)\}\).  Then
\(X_4=\{s_{a\,f_0(a)\,c\,f_1(c)}:a,c\in A\}\).  Two members with the same
\(x_1\)-value \(c\) both have \(y_1\)-value \(f_1(c)\), so \(=(x_1,y_1)\)
holds.  Each member satisfies \(\psi\) by the choice of \(f_0,f_1\).

\emph{Right to left.}  Let \(F_0,F_1\) be as above.  Fix \(c\in A\).  Since
\(A\neq\varnothing\) and the values of \(F_0\) and \(F_1\) are nonempty,
some member of \(X_4\) has \(x_1\)-value \(c\); by the dependence atom all
such members share one \(y_1\)-value, which we call \(f_1(c)\).  By the
Axiom of Choice pick \(f_0(a)\in F_0(s_a)\) for each \(a\in A\).  For any
\(a_0,a_1\in A\), the assignment \(s_{a_0f_0(a_0)a_1}\) lies in \(X_3\), and
for any \(d\in F_1(s_{a_0f_0(a_0)a_1})\) the assignment
\(s_{a_0f_0(a_0)a_1d}\) lies in \(X_4\); hence \(d=f_1(a_1)\).  So
\(s_{a_0f_0(a_0)a_1f_1(a_1)}\in X_4\) and it satisfies \(\psi\).  This is
the clause of \cref{def:p03-w03-p03-044-qh} for \(f_0,f_1\).
\end{proof}
\end{annexitem}
